\section{Numerical Results and Data Analysis}
\label{sec:data_analysis}

In this section, we present a comprehensive analysis of the numerical results obtained from our lattice QCD simulations. The primary objective is to investigate the restoration of chiral and $U(1)_A$ symmetries in $(N_f=2+1)$ flavor QCD, using Mobius Domain-Wall Fermions (MDWF). Our analysis is structured into two main parts:
\begin{itemize}
    \item Physical point simulations on a fixed $48^3 \times 16$ lattice across various inverse couplings $\beta$ (corresponding to different physical temperatures $T$).
    \item Fixed coupling simulations at $\beta=4.17$ with varying light quark masses and temporal lattice sizes to analyze the behavior in the chiral limit.
\end{itemize}

\subsection{Lattice Setup and Formulation}
For the gauge field, we employ the tree-level Symanzik-improved gauge action:
\begin{equation}
S_{\text{YM}}^{(L)}[U] = -\frac{\beta}{3} \sum_{x, \mu \neq \nu} \left\{ \frac{5}{6} \Re \text{Tr} \left[ \Pi_{\mu\nu}^{(1 \times 1)}(x) \right] - \frac{1}{12} \Re \text{Tr} \left[ \Pi_{\mu\nu}^{(1 \times 2)}(x) \right] \right\}
\label{eq:symanzik_action}
\end{equation}
where $\beta = 6/g^2$ is the inverse coupling, and $\Pi_{\mu\nu}^{(1 \times 1)}(x)$ and $\Pi_{\mu\nu}^{(1 \times 2)}(x)$ represent the plaquette ($1 \times 1$ loop) and the rectangular ($1 \times 2$ loop) in the $\mu$-$\nu$ plane, respectively. For the valence and sea quarks, we employ the M\"obius Domain-Wall Fermion (MDWF) formulation. The effective four-dimensional Dirac operator is given by:
\begin{equation}
\slashed{D}_{\mathrm{MDWF}}^{(4D)}(m) = \frac{1+m}{2} + \frac{1-m}{2} \gamma_5 \tanh \left( L_s \tanh^{-1}(H_M) \right)
\label{eq:mdwf_effective_analysis}
\end{equation}
where $m$ is the bare quark mass, $L_s$ is the size of the fifth dimension (we set $L_s = 16$ in our simulations), and the M\"obius kernel $H_M$ is:
\begin{equation}
H_M = \gamma_5 \frac{2 \slashed{D}_W}{2 + \slashed{D}_W}
\label{eq:mobius_hamiltonian_analysis}
\end{equation}
Here, $\slashed{D}_W$ is the standard massless 4D Wilson-Dirac operator with a negative mass parameter. The violation of the chiral Ward-Takahashi identity due to the finite fifth dimension $L_s$ is described by the contribution operator $\Delta_L$ we defined in the former section. In the low-energy effective field theory, this lattice chiral symmetry breaking operator can be expanded in terms of continuum operators as:
\begin{equation}
\overline{\psi} \Delta_L \psi \simeq m_{\mathrm{res}} \overline{\psi} \psi + c_1 \overline{\psi} \gamma_\mu (\partial_\mu - i A_\mu) \psi + i c_2 \overline{\psi} \sigma_{\mu\nu} F^{\mu\nu} \psi + c_4 \overline{\psi} (\partial_\mu - i A_\mu)^2 \psi + \cdots
\label{eq:midpoint_operator_expansion}
\end{equation}
where the leading contribution comes from the dimension-3 scalar density operator $\overline{\psi} \psi$. The coefficient of this term is the residual mass $m_{\mathrm{res}}$, which parameterizes the leading explicit chiral symmetry breaking. The precision of chiral symmetry on the lattice is controlled by this residual mass $m_{\mathrm{res}}$, which arises from the finite fifth dimension $L_s$.
In practice, we could use an effective definition of residual mass from analysis of one loop diagram:
$m_{\mathrm{res}} = \frac{1}{a} \frac{\langle 0 | J_{5q}^a | \pi^a \rangle}{\langle 0 | J_5^a | \pi^a \rangle}$, which is the correlation measurement result:
$$m_{\mathrm{res}} = R(t) = \frac{\sum_{\vec{x}} \langle J_{5q}(\vec{x}, t) P(\vec{0}, 0) \rangle}{\sum_{\vec{x}} \langle P(\vec{x}, 0) P(\vec{0}, 0) \rangle}$$.

We generate the $N_f=2+1$ gauge configurations using the Hybrid Monte Carlo (HMC) algorithm. Our simulations at the physical point are performed on a fixed $48^3 \times 16$ lattice. For each gauge coupling $\beta$, we collect several hundred configurations after thermalization to perform statistical measurements. The simulation parameters, including the physical temperatures, residual masses $m_{\rm res}$, bare quark masses ($am_s$, $am$), and the number of analyzed configurations $N_{\rm conf}$, are summarized in Tables \ref{tab:lattice_setup1}, \ref{tab:lattice_setup2}, and \ref{tab:lattice_confs}.  And our simulations at the fixed beta parameter 4.17 are performed on different scale lattice and light quark mass, also the different physical strange quark mass although the dimensionless parameter is the same, we print it at the table \ref{tab:lattice_confs}.

\begin{table}[H]
\centering
\renewcommand{\arraystretch}{1.2}
\small
\csvreader[no head,
  tabular={lcccccccc},
  table head=\hline\hline,
  late after first line=\\\hline,
  late after line=\\,
  late after last line=\\\hline\hline
]{pos/lattice_setup1.csv}{}{%
  \csvcoli & \csvcolii & \csvcoliii & \csvcoliv & \csvcolv & \csvcolvi & \csvcolvii & \csvcolviii & \csvcolix
}
\caption{Temperatures and corresponding residual masses for the physical point ensembles.}
\label{tab:lattice_setup1}
\end{table}

\begin{table}[H]
\centering
\renewcommand{\arraystretch}{1.2}
\small
\csvreader[no head,
  tabular={ccc},
  table head=\hline\hline,
  late after first line=\\\hline,
  late after line=\\,
  late after last line=\\\hline\hline
]{pos/lattice_setup2.csv}{}{%
  \csvcoli & \csvcolii & \csvcoliii
}
\caption{Bare quark masses for the strange and light quarks across simulated $\beta$ values.}
\label{tab:lattice_setup2}
\end{table}

\begin{table}[H]
\centering
\renewcommand{\arraystretch}{1.2}
\small
\csvreader[no head,
  tabular={lccccccc},
  table head=\hline\hline,
  late after first line=\\\hline,
  late after line=\\,
  late after last line=\\\hline\hline
]{pos/lattice_confs.csv}{}{%
  \csvcoli & \csvcolii & \csvcoliii & \csvcoliv & \csvcolv & \csvcolvi & \csvcolvii & \csvcolviii
}
\caption{Number of analyzed gauge configurations ($N_{\mathrm{conf}}$) and total measurement counts (using 16 sources per configuration) for each inverse coupling $\beta$ and temperature $T$.}
\label{tab:lattice_confs}
\end{table}

\begin{table}[H]
\centering
\renewcommand{\arraystretch}{1.2}
\caption{Simulation parameters at fixed $\beta = 4.17$ with Möbius Domain-Wall Fermions, including the lattice sizes, bare quark masses ($am_s$, $am_l$), and physical temperatures.}
\label{tab:fixed_beta_setup}
\begin{tabular}{ccccc}
\hline\hline
Lattice Size & $am_s$ & $am_l$ & $T$ (MeV) \\
\hline
$32^3 \times 12$ & 0.040 & 0.0020, 0.0035, 0.0070, 0.0120 & 220.00 \\
$32^3 \times 14$ & 0.040 & 0.0020, 0.0035, 0.0070, 0.0120 & 188.57 \\
$32^3 \times 16$ & 0.040 & 0.0020, 0.0035, 0.0070, 0.0120 & 165.00 \\
$36^3 \times 18$ & 0.040 & 0.0020, 0.0035, 0.0070, 0.0120 & 146.67 \\
$40^3 \times 16$ & 0.040 & 0.0020, 0.0035, 0.0070, 0.0120 & 165.00 \\
$48^3 \times 18$ & 0.040 & 0.0020, 0.0035, 0.0070, 0.0120 & 146.67 \\
\hline\hline
\end{tabular}
\end{table}


\subsection{Measurement and Analysis Techniques}
To compute the meson screening masses, we measure the spatial correlation functions $C(x)$ along each spatial axis. At each separation $x$, we compute the local effective mass $m_{\rm eff}(x)$ by numerically solving the ratio:
\begin{equation}
    \frac{C(x)}{C(x+1)} = \frac{\cosh\!\left(m_{\rm eff}(x)\bigl(x - \tfrac{L}{2}\bigr)\right)}{\cosh\!\left(m_{\rm eff}(x)\bigl(x + 1 - \tfrac{L}{2}\bigr)\right)}
\end{equation}

In practice, we use solver to solve:
\begin{equation}
    m_{\rm eff}(x) = \operatorname{arccosh}\left[\frac{C(x-1)+C(x+1)}{C(x)}\right]
\end{equation}

The resulting effective mass plateau is used to determine a fitting range $[x_{\rm min}, x_{\rm max}]$ to perform a correlated $\chi^2$ fit to the asymptotic form:
\begin{equation}
    C(x) = A\,\cosh\!\left(m_{\rm scr}\Bigl(x - \frac{L}{2}\Bigr)\right)
\end{equation}
from which the screening mass $m_{\rm scr}$ is extracted.

To reduce statistical noise, we use two averaging strategies:
\begin{itemize}
    \item Multi-source averaging: Multiple independent stochastic sources are placed at different lattice sites and their resulting correlators are averaged. In our work, we compare 16-source averaging (multi-source) with 1-source averaging (single-source) to demonstrate the efficiency of noise reduction.
    \item Directional averaging: Spatial correlators measured independently along the $x$-, $y$-, and $z$-directions are averaged, gaining a factor of $\sqrt{3}$ in statistical precision. The directional averaging components configured in our simulation are summarized in Table \ref{tab:averaging_components}.
\end{itemize}

\begin{table}[H]
\centering
\renewcommand{\arraystretch}{1.2}
\small
\setlength{\tabcolsep}{0.5pt}
\csvreader[no head,
  tabular={lcl},
  table head=\hline\hline,
  late after first line=\\\hline,
  late after line=\\,
  late after last line=\\\hline\hline
]{pos/averaging_components.csv}{}{%
  \csvcoli & \csvcolii & \csvcoliii
}
\caption{Summary of meson operators, components, and spatial directions averaged to perform directional averaging.}
\label{tab:averaging_components}
\end{table}

To handle statistical uncertainties and autocorrelations, we perform a Binning analysis to group consecutive trajectories into independent blocks.\cite{flyvbjergErrorEstimatesAverages1989} Given a primary dataset of gauge configurations $\{x_t\}$ for $t = 1, \dots, N_{\text{cfg}}$, we group them into $N_{\text{bin}}$ blocks of size $B = N_{\text{cfg}} / N_{\text{bin}}$. The binned block average $X_k$ for the $k$-th block is defined as:
\begin{equation}
X_k = \frac{1}{B} \sum_{j=1}^{B} x_{(k-1)B + j}, \quad k = 1, \dots, N_{\text{bin}}
\label{eq:binning_formula}
\end{equation}
This is followed by Jackknife resampling to calculate unbiased statistical errors. For an estimator $f$, the $i$-th Jackknife sample is computed by omitting the $i$-th bin: $\theta_i = f\left( \frac{1}{N_{\text{bin}} - 1} \sum_{j \neq i}^{N_{\text{bin}}} X_j \right)$. The Jackknife estimate of the statistical error is given by:
\begin{equation}
\sigma_{\theta} = \sqrt{ \frac{N_{\text{bin}} - 1}{N_{\text{bin}}} \sum_{i=1}^{N_{\text{bin}}} (\theta_i - \overline{\theta})^2 }
\label{eq:jackknife_error_formula}
\end{equation}
where $\overline{\theta} = \frac{1}{N_{\text{bin}}} \sum_{i=1}^{N_{\text{bin}}} \theta_i$ represents the average of the Jackknife samples. An example of the raw correlation data and the correspondingly extracted effective mass for the Vector channel at $\beta=4.15$ ($T=153$ MeV) are shown in Figure \ref{fig:corr_meff_example}.

\begin{figure}[H]
    \centering
    \begin{minipage}{0.48\textwidth}
        \centering
        \includegraphics[width=\linewidth]{pos/diff_clean_ps_plot.png}
        \caption{Correlation data of the Vec channel for $\beta=4.15$ ($T=153$ MeV).}
        \label{fig:corr_example}
    \end{minipage}\hfill
    \begin{minipage}{0.48\textwidth}
        \centering
        \includegraphics[width=\linewidth]{pos/outdia/b4.15/ratioVec.pdf}
        \caption{Effective mass of the Vec channel for $\beta=4.15$ ($T=153$ MeV).}
        \label{fig:meff_example}
    \end{minipage}
    \caption{Raw correlation function and the correspondingly extracted effective mass plateau for the Vector channel at $\beta=4.15$ ($T=153$ MeV).}
    \label{fig:corr_meff_example}
\end{figure}

To monitor the quality of our data fits and verify the stability of our statistical analysis, we check the $\chi^2$ value for each bin. To illustrate this step, Table \ref{tab:fit_sample_example} presents a representative sample of our fit results (specifically $\chi^2$, degrees of freedom $N_{\mathrm{df}}$ or dof, screening mass $a m_{\mathrm{screen}}$, and amplitude $a$) extracted from the Jackknife binning analysis.

\begin{table}[H]
\centering
\renewcommand{\arraystretch}{1.2}
\resizebox{\textwidth}{!}{%
\csvreader[no head,
  tabular={cccccccc},
  table head=\hline\hline,
  late after first line=\\\hline,
  late after line=\\,
  late after last line=\\\hline\hline
]{pos/fit_sample_example.csv}{}{%
  \csvcoli & \csvcolii & \csvcoliii & \csvcoliv & \csvcolv & \csvcolvi & \csvcolvii & \csvcolviii
}%
}
\caption{Sample fitting parameters containing $\chi^2$, degrees of freedom, screening mass, and amplitude extracted from the Jackknife binning analysis.}
\label{tab:fit_sample_example}
\end{table}

\subsection{Numerical Results and Symmetry Restoration}

\subsubsection*{Chiral Condensate and Crossover Phase Transition}
The chiral condensate $\langle \bar{\psi}\psi \rangle$ is the order parameter for spontaneous chiral symmetry breaking. However, the bare chiral condensate is divergent in the continuum limit due to additive and multiplicative renormalization factors. In our $N_f=2+1$ simulation, we define the renormalized chiral condensate by subtracting the divergent part values. The physical observable chiral condensation is defined as $\langle \bar{\psi}\psi \rangle = \frac{T}{V} \frac{\partial \ln \mathcal{Z}}{\partial m}$ for finite temperature and volume, and in our lattice case, the corresponding dimensionless observable is $a^3 \langle \bar{\psi}\psi \rangle = \frac{1}{N_s^3 N_\tau} \frac{N_f}{4} \langle \mathrm{Tr} \slashed{D}_{\mathrm{MDWF}}(m)^{-1} \rangle$. We use subtraction $\Delta_{l,s}(T) = Z^{-1}_m( \langle \bar{\psi}\psi_{l,T} \rangle - \frac{m_l+ m_{\rm res}}{m_s+m_{\rm res}} \langle \bar{\psi}\psi_{s,T} \rangle)$ to remove the divergent strange quark condensation effect.

In Figure \ref{fig:condensates}, we show the bare and renormalized chiral condensates as functions of temperature. The bare condensate shows a gradual decrease with temperature. The renormalized chiral condensate provides a much clearer picture of the thermal transition, where we observe a smooth crossover behavior with the condensate melting. The transition occurs in the pseudo-critical temperature region $T_c \approx 155\text{--}160$ MeV. This crossover signals the transition from the low-temperature chiral-broken phase to the high-temperature chiral-restored phase of QCD. The measured values for the bare light quark condensate, strange quark condensate, subtracted condensate, and the mass renormalization factor $Z_m$ are summarized in Table \ref{tab:chiral_condensate_data}.

\begin{table}[H]
\centering
\renewcommand{\arraystretch}{1.2}
\caption{Bare light quark condensate $a^3\langle \bar{\psi}\psi_l \rangle$, strange quark condensate $a^3\langle \bar{\psi}\psi_s \rangle$, subtracted condensate $\Delta_{l,s}$, and mass renormalization factor $Z_m$ for different inverse couplings $\beta$.}
\label{tab:chiral_condensate_data}
\begin{tabular}{ccccc}
\hline\hline
$\beta$ &
\begin{tabular}{c} $a^3\langle \bar{\psi}\psi_l \rangle$ \\ (Light Quark) \end{tabular} &
\begin{tabular}{c} $a^3\langle \bar{\psi}\psi_s \rangle$ \\ (Strange Quark) \end{tabular} &
\begin{tabular}{c} $\Delta_{l,s}$ \\ (Subtracted) \end{tabular} &
$Z_m$ \\
\hline
4.13  & \begin{tabular}{c} $0.00166880$ \\ $\pm 0.00000634$ \end{tabular} & \begin{tabular}{c} $0.01942524$ \\ $\pm 0.00000135$ \end{tabular} & \begin{tabular}{c} $0.00106082$ \\ $\pm 0.00000674$ \end{tabular} & $0.937703$ \\
4.15  & \begin{tabular}{c} $0.00135913$ \\ $\pm 0.00000644$ \end{tabular} & \begin{tabular}{c} $0.01808069$ \\ $\pm 0.00000128$ \end{tabular} & \begin{tabular}{c} $0.00077109$ \\ $\pm 0.00000675$ \end{tabular} & $0.952390$ \\
4.17  & \begin{tabular}{c} $0.00107396$ \\ $\pm 0.00000642$ \end{tabular} & \begin{tabular}{c} $0.01686886$ \\ $\pm 0.00000127$ \end{tabular} & \begin{tabular}{c} $0.00050728$ \\ $\pm 0.00000662$ \end{tabular} & $0.966247$ \\
4.18  & \begin{tabular}{c} $0.00097061$ \\ $\pm 0.00000642$ \end{tabular} & \begin{tabular}{c} $0.01629887$ \\ $\pm 0.00000134$ \end{tabular} & \begin{tabular}{c} $0.00041649$ \\ $\pm 0.00000659$ \end{tabular} & $0.972899$ \\
4.20  & \begin{tabular}{c} $0.00076708$ \\ $\pm 0.00000554$ \end{tabular} & \begin{tabular}{c} $0.01524584$ \\ $\pm 0.00000127$ \end{tabular} & \begin{tabular}{c} $0.00023896$ \\ $\pm 0.00000561$ \end{tabular} & $0.985710$ \\
4.23  & \begin{tabular}{c} $0.00058018$ \\ $\pm 0.00000383$ \end{tabular} & \begin{tabular}{c} $0.01384239$ \\ $\pm 0.00000120$ \end{tabular} & \begin{tabular}{c} $0.00009289$ \\ $\pm 0.00000381$ \end{tabular} & $1.003850$ \\
4.30  & \begin{tabular}{c} $0.00041210$ \\ $\pm 0.00000139$ \end{tabular} & \begin{tabular}{c} $0.01124911$ \\ $\pm 0.00000090$ \end{tabular} & \begin{tabular}{c} $0.00000820$ \\ $\pm 0.00000133$ \end{tabular} & $1.042240$ \\
4.405 & \begin{tabular}{c} $0.00031129$ \\ $\pm 0.00000003$ \end{tabular} & \begin{tabular}{c} $0.00856161$ \\ $\pm 0.00000027$ \end{tabular} & \begin{tabular}{c} $0.00000040$ \\ $\pm 0.00000002$ \end{tabular} & $1.092740$ \\
\hline\hline
\end{tabular}
\end{table}

\begin{figure}[H]
    \centering
    \begin{minipage}{0.48\textwidth}
        \centering
        \includegraphics[width=\linewidth]{pos/conden_plot.png}
    \end{minipage}\hfill
    \begin{minipage}{0.48\textwidth}
        \centering
        \includegraphics[width=\linewidth]{pos/conden_re_plot.png}
    \end{minipage}
    \caption{Bare chiral condensate (left) and Renormalized chiral condensate (right) as a function of temperature.}
    \label{fig:condensates}
\end{figure}

Furthermore, we measure the chiral susceptibility, which is defined as the derivative of the chiral condensate with respect to the light quark mass. For our purpose, we show disconnected susceptibility result correlated to the topological effect, defined as $a^2 \chi_{dis} = \frac{1}{N_s^3 N_\tau} \left( \frac{N_f}{4} \right)^2 \left( \langle (\mathrm{Tr} \slashed{D}_{\mathrm{MDWF}}(m)^{-1})^2 \rangle -\right.$ $\left.\langle \mathrm{Tr} \slashed{D}_{\mathrm{MDWF}}(m)^{-1} \rangle^2 \right)$. The chiral susceptibility at the physical point is plotted in Figure \ref{fig:chiral_susceptibility} as a function of temperature. We observe a clear peak in the temperature range $T \approx 155\text{--}165$ MeV, which defines the pseudo-critical transition temperature $T_c$. This peak is consistent with the chiral crossover transition at the physical point and aligns with the melting of the renormalized chiral condensate. The numerical data for the susceptibility measurements are given in Table \ref{tab:chiral_susceptibility_data}.

\begin{table}[H]
\centering
\renewcommand{\arraystretch}{1.2}
\caption{Disconnected chiral susceptibility $\chi_{\text{dis}}$ at the physical point across different couplings $\beta$ and temperatures $T$, showing both the unscaled lattice value $a^2 \chi_{\text{dis}}$ and the scaled physical value.}
\label{tab:chiral_susceptibility_data}
\begin{tabular}{cccc}
\hline\hline
$\beta$ &
\begin{tabular}{c} Temperature \\ $T$ [MeV] \end{tabular} &
\begin{tabular}{c} Unscaled \\ $a^2 \chi_{\text{dis}}$ \end{tabular} &
\begin{tabular}{c} Scaled \\ $\chi_{\text{dis}}$ \end{tabular} \\
\hline
4.13  & 138.0 & \begin{tabular}{c} $7.650$ \\ $\pm 0.266$ \end{tabular} $\times 10^{-8}$  & $0.1354 \pm 0.0047$ \\
4.15  & 145.0 & \begin{tabular}{c} $8.106$ \\ $\pm 0.278$ \end{tabular} $\times 10^{-8}$  & $0.1434 \pm 0.0049$ \\
4.17  & 153.0 & \begin{tabular}{c} $8.061$ \\ $\pm 0.242$ \end{tabular} $\times 10^{-8}$  & $0.1426 \pm 0.0043$ \\
4.18  & 157.0 & \begin{tabular}{c} $8.191$ \\ $\pm 0.289$ \end{tabular} $\times 10^{-8}$  & $0.1449 \pm 0.0051$ \\
4.20  & 164.5 & \begin{tabular}{c} $6.245$ \\ $\pm 0.226$ \end{tabular} $\times 10^{-8}$  & $0.1105 \pm 0.0040$ \\
4.23  & 176.0 & \begin{tabular}{c} $2.901$ \\ $\pm 0.164$ \end{tabular} $\times 10^{-8}$  & $0.0513 \pm 0.0029$ \\
4.30  & 202.0 & \begin{tabular}{c} $3.871$ \\ $\pm 0.809$ \end{tabular} $\times 10^{-9}$  & $0.0068 \pm 0.0014$ \\
4.405 & 241.6 & \begin{tabular}{c} $3.271$ \\ $\pm 0.641$ \end{tabular} $\times 10^{-12}$ & \begin{tabular}{c} $5.79$ \\ $\pm 1.13$ \end{tabular} $\times 10^{-6}$ \\
\hline\hline
\end{tabular}
\end{table}

\begin{figure}[H]
    \centering
    \includegraphics[width=0.6\linewidth]{pos/pbpchisce.png}
    \caption{Chiral susceptibility as a function of temperature at the physical point, showing a peak near the pseudo-critical temperature region.}
    \label{fig:chiral_susceptibility}
\end{figure}

\subsubsection*{Screening Masses and Symmetry Restoration at the Physical Point}
We analyze the screening masses extracted from both multi-source and single-source measurements across all meson channels on the $48^3 \times 16$ lattice. The multi-source measurements (using 16 independent sources per configuration) and single-source measurements (using 1 source per configuration) allow us to compare the effectiveness of noise reduction. The numerical values of the screening masses for multi-source measurements are detailed in Table \ref{tab:screening_masses_data_multi}, and those for single-source measurements are detailed in Table \ref{tab:screening_masses_data_single}. Furthermore, the corresponding symmetry breaking mass differences for partner channels are summarized in Table \ref{tab:symmetry_breaking_data_multi} (multi-source) and Table \ref{tab:symmetry_breaking_data_single} (single-source).

A visual comparison of the screening masses between multi-source and single-source calculations is shown in Figure \ref{fig:all_simulation_results_comparison}. In the single-source results (right panel), the data points are highly scattered with extremely large error bars, especially for the Scalar ($S$), Tensor ($Tt$), and Axial-Tensor ($Xt$) channels, where the signal-to-noise ratio degrades rapidly at large spatial separations. In contrast, the multi-source results (left panel) show a much smoother temperature dependence and significantly smaller statistical errors (reduced by a factor of 4 to 8).

Similarly, Figure \ref{fig:symmetry_breaking_comparison} compares the symmetry breaking mass differences between multi-source and single-source configurations. While the single-source measurements are too noisy to draw definitive physical conclusions, the multi-source measurements clearly reveal the systematic restoration of symmetries.

\begin{table}[H]
\centering
\renewcommand{\arraystretch}{1.3}
\resizebox{\textwidth}{!}{%
\csvreader[no head,
  tabular={cccccccc},
  table head=\hline\hline,
  late after first line=\\\hline,
  late after line=\\,
  late after last line=\\\hline\hline
]{pos/table.csv}{}{%
\csvcolii & \csvcoliii & \csvcoliv & \csvcolv & \csvcolvi & \csvcolvii & \csvcolviii & \csvcolix }%
}
\caption{Screening masses (in MeV) for various meson channels at different temperatures at the physical point, obtained from multi-source measurements (16 sources).}
\label{tab:screening_masses_data_multi}
\end{table}

\begin{table}[H]
\centering
\renewcommand{\arraystretch}{1.3}
\resizebox{\textwidth}{!}{%
\csvreader[no head,
  tabular={ccccccc},
  table head=\hline\hline,
  late after first line=\\\hline,
  late after line=\\,
  late after last line=\\\hline\hline
]{pos/table.csv}{}{%
  \csvcolii & \csvcolx & \csvcolxi & \csvcolxii & \csvcolxiii & \csvcolxiv & \csvcolxv
}%
}
\caption{Screening masses (in MeV) for various meson channels at different temperatures at the physical point, obtained from single-source measurements (1 source).}
\label{tab:screening_masses_data_single}
\end{table}

\begin{table}[H]
\centering
\renewcommand{\arraystretch}{1.3}
\resizebox{\textwidth}{!}{%
\csvreader[no head,
  tabular={ccccc},
  table head=\hline\hline,
  late after first line=\\\hline,
  late after line=\\,
  late after last line=\\\hline\hline
]{pos/table_sb.csv}{}{%
  \csvcolii & \csvcolvi & \csvcolviii & \csvcolx & \csvcolxii
}%
}
\caption{Symmetry breaking mass differences (in MeV) between partner channels at different temperatures at the physical point, obtained from multi-source measurements (16 sources).}
\label{tab:symmetry_breaking_data_multi}
\end{table}

\begin{table}[H]
\centering
\renewcommand{\arraystretch}{1.3}
\resizebox{\textwidth}{!}{%
\csvreader[no head,
  tabular={ccccc},
  table head=\hline\hline,
  late after first line=\\\hline,
  late after line=\\,
  late after last line=\\\hline\hline
]{pos/table_sb.csv}{}{%
  \csvcolii & \csvcolvii & \csvcolix & \csvcolxi & \csvcolxiii
}%
}
\caption{Symmetry breaking mass differences (in MeV) between partner channels at different temperatures at the physical point, obtained from single-source measurements (1 source).}
\label{tab:symmetry_breaking_data_single}
\end{table}

\begin{figure}[H]
    \centering
    \begin{subfigure}{0.48\textwidth}
        \centering
        \includegraphics[width=\linewidth]{pos/simulateresult/simulation_results.pdf}
        \caption{Multi-source measurements (16 sources)}
        \label{fig:all_simulation_results_multi}
    \end{subfigure}\hfill
    \begin{subfigure}{0.48\textwidth}
        \centering
        \includegraphics[width=\linewidth]{pos/simulateresult-singlesrc/simulation_results.pdf}
        \caption{Single-source measurements}
        \label{fig:all_simulation_results_single}
    \end{subfigure}
    \caption{Comparison of meson screening masses across all channels as a function of temperature at the physical point between multi-source (left) and single-source (right) measurements.}
    \label{fig:all_simulation_results_comparison}
\end{figure}

\begin{figure}[H]
    \centering
    \begin{subfigure}{0.48\textwidth}
        \centering
        \includegraphics[width=\linewidth]{pos/simulateresult/symmetry_breaking.pdf}
        \caption{Multi-source measurements (16 sources)}
        \label{fig:symmetry_breaking_multi}
    \end{subfigure}\hfill
    \begin{subfigure}{0.48\textwidth}
        \centering
        \includegraphics[width=\linewidth]{pos/simulateresult-singlesrc/symmetry_breaking.pdf}
        \caption{Single-source measurements}
        \label{fig:symmetry_breaking_single}
    \end{subfigure}
    \caption{Comparison of symmetry breaking mass differences as a function of temperature at the physical point between multi-source (left) and single-source (right) measurements. The grey band indicates the crossover region $T_c \approx 158$ MeV.}
    \label{fig:symmetry_breaking_comparison}
\end{figure}

Let us discuss the physical features of the high-quality multi-source data:
\begin{itemize}
    \item Low-Temperature Regime ($T < 155$ MeV): In the broken phase, the screening masses of the partner channels differ significantly. In particular, the pseudoscalar (pion) screening mass is light ($m_{PS} = 162.29(1.75)$ MeV at $T=138.23$ MeV). Meanwhile, the scalar channel mass is $m_S = 23.06(215.42)$ MeV with a very large uncertainty due to the low signal-to-noise ratio in the scalar correlation function (a known challenge in lattice QCD simulations of the scalar $a_0$ channel).
    \item High-Temperature Regime ($T > 175$ MeV): Above the critical crossover temperature, we observe a clear trend supporting the restoration of both chiral and axial symmetries.
    \begin{itemize}
        \item Chiral $SU(2)_L \times SU(2)_R$ restoration: The vector and axial-vector channels display a high degree of degeneracy. At $T = 204.52$ MeV, the mass splitting is $m_{Vec} - m_{AV} = 1.07(11.85)$ MeV, which is consistent with zero. The details of this restoration are shown in Figure \ref{fig:symmetry_breaking_Chiral_AV_Vec}. At $T=0$, the difference between the $a_1$ ($1230$ MeV) and $\rho$ meson ($775$ MeV) is around $455$ MeV, showing that while the symmetry breaking is very strong in the vacuum, it is consistent with being restored at high temperatures within our current statistical uncertainties.
        \item Axial $U(1)_A$ restoration: The scalar and pseudoscalar channels exhibit degeneracy, with a splitting of $m_S - m_{PS} = -0.47(130.04)$ MeV at $T = 204.52$ MeV (shown in Figure \ref{fig:symmetry_breaking_U1A_S_PS}). The tensor and axial-tensor channels also show degeneracy with a splitting of $m_{Xt} - m_{Tt} = 6.28(25.34)$ MeV (shown in Figure \ref{fig:symmetry_breaking_U1A_Xt_Tt}). This is consistent with the expected suppression of topological fluctuations (instantons) in this regime, pointing toward an effective restoration of the anomalous $U(1)_A$ symmetry.
        \item Emergent Chiral-Spin Symmetry ($SU(2)_{CS}$): The chiral-spin symmetry is investigated via the mass difference between the AV and Tt channels, shown in Figure \ref{fig:symmetry_breaking_Xt_AV}. Unlike chiral and $U(1)_A$ symmetries, the mass splitting remains large up to $T = 204.52$ MeV ($m_{AV} - m_{Tt} = -163.61(19.14)$ MeV). This suggests that the chiral-spin symmetry remains broken in the analyzed temperature range and chromomagnetic interactions remain active.
    \end{itemize}
\end{itemize}

\begin{figure}[H]
    \centering
    \includegraphics[width=0.65\linewidth]{pos/simulateresult/symmetry_breaking_Chiral_AV_Vec.pdf}
    \caption{Symmetry breaking mass difference $m_{Vec}-m_{AV}$ as a function of temperature. The grey band indicates the crossover region $T_c \approx 158$ MeV.}
    \label{fig:symmetry_breaking_Chiral_AV_Vec}
\end{figure}

\begin{figure}[H]
    \centering
    \includegraphics[width=0.65\linewidth]{pos/simulateresult/symmetry_breaking_U1A_S_PS.pdf}
    \caption{Symmetry breaking mass difference $m_{S}-m_{PS}$ as a function of temperature. The grey band indicates the crossover region $T_c \approx 165(3)$ MeV.}
    \label{fig:symmetry_breaking_U1A_S_PS}
\end{figure}

\begin{figure}[H]
    \centering
    \includegraphics[width=0.65\linewidth]{pos/simulateresult/symmetry_breaking_U1A_Xt_Tt.pdf}
    \caption{Symmetry breaking mass difference $m_{Xt}-m_{Tt}$ as a function of temperature. The grey band indicates the crossover region $T_c \approx 165(3)$ MeV.}
    \label{fig:symmetry_breaking_U1A_Xt_Tt}
\end{figure}

\begin{figure}[H]
    \centering
    \includegraphics[width=0.65\linewidth]{pos/simulateresult/symmetry_breaking_Xt_AV.pdf}
    \caption{Symmetry breaking mass difference between $AV$ and $T_t$ channels testing the emergent $SU(2)_{CS}$ chiral-spin symmetry.}
    \label{fig:symmetry_breaking_Xt_AV}
\end{figure}

\subsubsection*{Symmetry Restoration at Fixed $\beta=4.17$}
To study the restoration mechanism in the chiral limit, we perform simulations at a fixed coupling $\beta=4.17$, varying the bare light quark mass $am_l \in \{0.0020, 0.0035, 0.0070, 0.0120\}$ and the temporal lattice size $N_t$. The corresponding lattice setups, bare quark masses, and physical temperatures ranging from $147$ MeV to $220$ MeV are summarized in Table \ref{tab:fixed_beta_setup}. The raw lattice mass differences ($a\Delta M$) and their physical conversions ($\Delta M$ in MeV) for the four symmetry partner combinations ($SU(2)_L \times SU(2)_R$, axial $U(1)_A$, axial $U(1)$, and chiral-spin $SU(2)_{CS}$) are tabulated in Tables \ref{tab:su2_l_su2_r_asym_data}, \ref{tab:u1_a_tx_asym_data}, \ref{tab:u1_s_ps_asym_data}, and \ref{tab:su2_cs_xa_asym_data}, respectively.

\begin{table}[H]
\centering
\renewcommand{\arraystretch}{1.2}
\caption{Symmetry breaking mass differences for Chiral $SU(2)_L \times SU(2)_R$ ($V-A$) in lattice units ($a\Delta M$) and physical units ($\Delta M$ in MeV) at fixed $\beta=4.17$.}
\label{tab:su2_l_su2_r_asym_data}
\begin{tabular}{cccc}
\hline\hline
$T$ (MeV) & $am_l$ & $a\Delta M$ & $\Delta M$ (MeV) \\
\hline
146.67 & 0.0020 & $0.01391 \pm 0.01590$ & $34.12 \pm 38.99$ \\
146.67 & 0.0035 & $0.04106 \pm 0.01902$ & $100.73 \pm 46.66$ \\
146.67 & 0.0070 & $0.08210 \pm 0.01999$ & $201.40 \pm 49.03$ \\
146.67 & 0.0120 & $0.08687 \pm 0.01781$ & $213.08 \pm 43.69$ \\
\hline
165.00 & 0.0020 & $0.01324 \pm 0.00773$ & $32.48 \pm 18.96$ \\
165.00 & 0.0035 & $0.03555 \pm 0.01326$ & $87.21 \pm 32.53$ \\
165.00 & 0.0070 & $0.04971 \pm 0.01313$ & $121.95 \pm 32.20$ \\
165.00 & 0.0120 & $0.04978 \pm 0.01349$ & $122.11 \pm 33.08$ \\
\hline
188.57 & 0.0020 & $0.00493 \pm 0.00542$ & $12.10 \pm 13.30$ \\
188.57 & 0.0035 & $0.01135 \pm 0.00752$ & $27.84 \pm 18.44$ \\
188.57 & 0.0070 & $0.04158 \pm 0.01142$ & $102.00 \pm 28.01$ \\
188.57 & 0.0120 & $0.04493 \pm 0.00947$ & $110.20 \pm 23.24$ \\
\hline
220.00 & 0.0020 & $-0.00045 \pm 0.00043$ & $-1.10 \pm 1.07$ \\
220.00 & 0.0035 & $0.00038 \pm 0.00198$ & $0.94 \pm 4.87$ \\
220.00 & 0.0070 & $0.00941 \pm 0.00462$ & $23.09 \pm 11.34$ \\
220.00 & 0.0120 & $0.01226 \pm 0.01626$ & $30.08 \pm 39.88$ \\
\hline\hline
\end{tabular}
\end{table}

\begin{table}[H]
\centering
\renewcommand{\arraystretch}{1.2}
\caption{Symmetry breaking mass differences for axial $U(1)_A$ ($T-X$) in lattice units ($a\Delta M$) and physical units ($\Delta M$ in MeV) at fixed $\beta=4.17$.}
\label{tab:u1_a_tx_asym_data}
\begin{tabular}{cccc}
\hline\hline
$T$ (MeV) & $am_l$ & $a\Delta M$ & $\Delta M$ (MeV) \\
\hline
146.67 & 0.0020 & $0.03342 \pm 0.02357$ & $81.98 \pm 57.82$ \\
146.67 & 0.0035 & $0.01866 \pm 0.02261$ & $45.78 \pm 55.47$ \\
146.67 & 0.0070 & $0.00059 \pm 0.02134$ & $1.46 \pm 52.35$ \\
146.67 & 0.0120 & $0.09245 \pm 0.01842$ & $226.77 \pm 45.19$ \\
\hline
165.00 & 0.0020 & $0.00573 \pm 0.00884$ & $14.06 \pm 21.68$ \\
165.00 & 0.0035 & $0.01639 \pm 0.00927$ & $40.20 \pm 22.74$ \\
165.00 & 0.0070 & $0.00487 \pm 0.01488$ & $11.94 \pm 36.50$ \\
165.00 & 0.0120 & $0.07532 \pm 0.01363$ & $184.76 \pm 33.44$ \\
\hline
188.57 & 0.0020 & $0.00386 \pm 0.00593$ & $9.48 \pm 14.55$ \\
188.57 & 0.0035 & $0.00617 \pm 0.00873$ & $15.13 \pm 21.42$ \\
188.57 & 0.0070 & $0.00849 \pm 0.00946$ & $20.84 \pm 23.21$ \\
188.57 & 0.0120 & $0.03312 \pm 0.00921$ & $81.24 \pm 22.58$ \\
\hline
220.00 & 0.0020 & $0.00399 \pm 0.00425$ & $9.78 \pm 10.43$ \\
220.00 & 0.0035 & $-0.00304 \pm 0.00222$ & $-7.46 \pm 5.44$ \\
220.00 & 0.0070 & $0.00188 \pm 0.00829$ & $4.61 \pm 20.32$ \\
220.00 & 0.0120 & $0.00628 \pm 0.00841$ & $15.41 \pm 20.62$ \\
\hline\hline
\end{tabular}
\end{table}

\begin{table}[H]
\centering
\renewcommand{\arraystretch}{1.2}
\caption{Symmetry breaking mass differences for axial $U(1)$ ($S-PS$) in lattice units ($a\Delta M$) and physical units ($\Delta M$ in MeV) at fixed $\beta=4.17$.}
\label{tab:u1_s_ps_asym_data}
\begin{tabular}{cccc}
\hline\hline
$T$ (MeV) & $am_l$ & $a\Delta M$ & $\Delta M$ (MeV) \\
\hline
146.67 & 0.0020 & $-0.01193 \pm 0.05191$ & $-29.26 \pm 127.33$ \\
146.67 & 0.0035 & $0.00490 \pm 0.01129$ & $12.02 \pm 27.70$ \\
146.67 & 0.0070 & $-0.00380 \pm 0.01841$ & $-9.33 \pm 45.15$ \\
146.67 & 0.0120 & $0.03826 \pm 0.01617$ & $93.85 \pm 39.66$ \\
\hline
165.00 & 0.0020 & $-0.00757 \pm 0.00923$ & $-18.58 \pm 22.64$ \\
165.00 & 0.0035 & $0.01764 \pm 0.00714$ & $43.28 \pm 17.52$ \\
165.00 & 0.0070 & $0.01706 \pm 0.00969$ & $41.85 \pm 23.76$ \\
165.00 & 0.0120 & $0.02987 \pm 0.01326$ & $73.28 \pm 32.52$ \\
\hline
188.57 & 0.0020 & $0.00333 \pm 0.00738$ & $8.17 \pm 18.11$ \\
188.57 & 0.0035 & $0.00065 \pm 0.00975$ & $1.60 \pm 23.93$ \\
188.57 & 0.0070 & $0.01230 \pm 0.00927$ & $30.17 \pm 22.73$ \\
188.57 & 0.0120 & $0.02718 \pm 0.01466$ & $66.68 \pm 35.96$ \\
\hline
220.00 & 0.0020 & $0.00309 \pm 0.00249$ & $7.58 \pm 6.11$ \\
220.00 & 0.0035 & $0.00185 \pm 0.00576$ & $4.54 \pm 14.14$ \\
220.00 & 0.0070 & $0.01335 \pm 0.00636$ & $32.74 \pm 15.59$ \\
220.00 & 0.0120 & $0.02411 \pm 0.00872$ & $59.14 \pm 21.38$ \\
\hline\hline
\end{tabular}
\end{table}

\begin{table}[H]
\centering
\renewcommand{\arraystretch}{1.2}
\caption{Symmetry breaking mass differences for chiral-spin $SU(2)_{CS}$ ($X-A$) in lattice units ($a\Delta M$) and physical units ($\Delta M$ in MeV) at fixed $\beta=4.17$.}
\label{tab:su2_cs_xa_asym_data}
\begin{tabular}{cccc}
\hline\hline
$T$ (MeV) & $am_l$ & $a\Delta M$ & $\Delta M$ (MeV) \\
\hline
146.67 & 0.0020 & $0.06088 \pm 0.04721$ & $149.33 \pm 115.82$ \\
146.67 & 0.0035 & $0.00606 \pm 0.03593$ & $14.87 \pm 88.13$ \\
146.67 & 0.0070 & $-0.02194 \pm 0.02687$ & $-53.81 \pm 65.91$ \\
146.67 & 0.0120 & $0.01479 \pm 0.02326$ & $36.27 \pm 57.05$ \\
\hline
165.00 & 0.0020 & $0.02279 \pm 0.02163$ & $55.91 \pm 53.05$ \\
165.00 & 0.0035 & $0.02086 \pm 0.02149$ & $51.17 \pm 52.71$ \\
165.00 & 0.0070 & $0.00454 \pm 0.01735$ & $11.14 \pm 42.56$ \\
165.00 & 0.0120 & $0.04310 \pm 0.01594$ & $105.73 \pm 39.09$ \\
\hline
188.57 & 0.0020 & $0.07138 \pm 0.02169$ & $175.09 \pm 53.21$ \\
188.57 & 0.0035 & $0.10993 \pm 0.01914$ & $269.65 \pm 46.94$ \\
188.57 & 0.0070 & $0.05772 \pm 0.02049$ & $141.59 \pm 50.27$ \\
188.57 & 0.0120 & $0.04827 \pm 0.01503$ & $118.42 \pm 36.87$ \\
\hline
220.00 & 0.0020 & $0.05405 \pm 0.02200$ & $132.58 \pm 53.96$ \\
220.00 & 0.0035 & $0.03899 \pm 0.03250$ & $95.64 \pm 79.72$ \\
220.00 & 0.0070 & $0.05294 \pm 0.02829$ & $129.86 \pm 69.41$ \\
220.00 & 0.0120 & $0.01031 \pm 0.01859$ & $25.29 \pm 45.60$ \\
\hline\hline
\end{tabular}
\end{table}

The temperature dependence of the screening masses for different light quark masses is shown in Figure \ref{fig:mass_fixed_beta}. As the temperature increases, partner channels in the PS-S and V-AV sectors become degenerate for all simulated light quark masses.

\begin{figure}[H]
    \centering
    \begin{subfigure}{0.48\textwidth}
        \centering
        \includegraphics[width=\linewidth]{pos/ccor/massvtem_0.0020.pdf}
        \caption{$am_l = 0.0020$}
        \label{fig:mass_fixed_beta_0020}
    \end{subfigure}\hfill
    \begin{subfigure}{0.48\textwidth}
        \centering
        \includegraphics[width=\linewidth]{pos/ccor/massvtem_0.0035.pdf}
        \caption{$am_l = 0.0035$}
        \label{fig:mass_fixed_beta_0035}
    \end{subfigure}

    \vspace{1em}

    \begin{subfigure}{0.48\textwidth}
        \centering
        \includegraphics[width=\linewidth]{pos/ccor/massvtem_0.0070.pdf}
        \caption{$am_l = 0.0070$}
        \label{fig:mass_fixed_beta_0070}
    \end{subfigure}\hfill
    \begin{subfigure}{0.48\textwidth}
        \centering
        \includegraphics[width=\linewidth]{pos/ccor/massvtem_0.0120.pdf}
        \caption{$am_l = 0.0120$}
        \label{fig:mass_fixed_beta_0120}
    \end{subfigure}
    \caption{Screening masses as a function of temperature for fixed $\beta=4.17$ and different light quark masses $am_l \in \{0.0020, 0.0035, 0.0070, 0.0120\}$.}
    \label{fig:mass_fixed_beta}
\end{figure}

The asymmetry observables for the different symmetry groups are plotted in Figure \ref{fig:asym_fixed_beta}. We observe that all symmetries tend towards restoration above $T_c$. The $U(1)_A$ and chiral $SU(2)_L \times SU(2)_R$ symmetries show an almost simultaneous restoration as we approach the chiral limit.

\begin{figure}[H]
    \centering
    \begin{minipage}{0.48\textwidth}
        \centering
        \includegraphics[width=\linewidth]{pos/ccor/dataSU(2)XSU(2) asym.pdf}
        \caption{$SU(2)_L \times SU(2)_R$ asymmetry (V-AV).}
    \end{minipage}\hfill
    \begin{minipage}{0.48\textwidth}
        \centering
        \includegraphics[width=\linewidth]{pos/ccor/dataU(1)A asym.pdf}
        \caption{$U(1)_A$ asymmetry (T-X).}
    \end{minipage}
    \vspace{1em}
    \begin{minipage}{0.48\textwidth}
        \centering
        \includegraphics[width=\linewidth]{pos/ccor/dataU(1) asym.pdf}
        \caption{$U(1)_A$ asymmetry (S-PS).}
    \end{minipage}\hfill
    \begin{minipage}{0.48\textwidth}
        \centering
        \includegraphics[width=\linewidth]{pos/ccor/dataSU(2)spinxchiral asym.pdf}
        \caption{$SU(2)_{CS}$ asymmetry (X-AV).}
    \end{minipage}
    \caption{Asymmetry observables as a function of temperature for various light quark masses at $\beta=4.17$.}
    \label{fig:asym_fixed_beta}
\end{figure}

Finally, to verify the restoration in the chiral limit, we analyze the mass differences at a high-temperature snapshot ($T \approx 220$ MeV) as a function of the bare light quark mass $am_l$. The data is shown in Figure \ref{fig:delta_mass_m}.
\begin{figure}[H]
    \centering
    \includegraphics[width=0.6\linewidth]{pos/ccor/deltamassvsm.pdf}
    \caption{Mass differences between mesonic partners as a function of $am_l$ at fixed $\beta=4.17$ and $T \approx 220$ MeV.}
    \label{fig:delta_mass_m}
\end{figure}
As the quark mass $am_l \to 0$, the mass differences between chiral partners ($m_{Vec}-m_{AV}$ and $m_{S}-m_{PS}$) extrapolate toward zero, supporting the scenario of a simultaneous restoration of chiral and $U(1)_A$ symmetries in the chiral limit.



\subsubsection*{Summary and Conclusion}
In this thesis, we have systematically investigated the chiral phase transition and the hierarchical restoration of symmetries in $N_{f}=2+1$ flavor QCD. The precise fate of the anomalous $U(1)_{A}$ symmetry at high temperatures has remained a highly debated topic, largely because conventional numerical simulations utilizing staggered fermions (such as HISQ) suffer from residual taste-symmetry breaking at finite lattice spacing. To overcome this theoretical bottleneck and avoid distorting the critical scaling behavior, we employed the Möbius Domain-Wall Fermion (MDWF) formalism. By strictly satisfying the Ginsparg-Wilson relation, our MDWF setup guarantees an exact chiral symmetry on the lattice, allowing us to unambiguously separate genuine physical symmetry restoration from lattice artifacts.Our numerical analysis of the renormalized chiral condensate and disconnected chiral susceptibility firmly identifies a smooth crossover transition with a pseudo-critical temperature $T_{c}$ around 156-160 MeV. By extracting and tracking the spatial screening masses across different mesonic parity-partner channels, we have identified the following phenomenological features of the high-temperature phase:
\begin{itemize}
    \item A smooth crossover transition in the pseudo-critical temperature region $T_c \approx 156\text{--}160$ MeV, characterized by the gradual melting of the chiral condensate.
    \item A clear trend toward the restoration of $SU(2)_L \times SU(2)_R$ chiral flavor symmetry at high temperatures, indicated by the degeneracy of vector and axial-vector screening masses within statistical uncertainties.
    \item Indications of an effective restoration of the axial $U(1)_A$ symmetry, as suggested by the degeneracy of the scalar-pseudoscalar and tensor-axial-tensor channels, consistent with the suppression of topological fluctuations above $T_c$.
    \item Extrapolations to the chiral limit at $T \approx 220$ MeV support the scenario of a simultaneous restoration of both symmetries.
    \item Chiral-spin symmetry $SU(2)_{CS}$ remains broken in the analyzed temperature range, indicating that chromomagnetic interactions remain active.
\end{itemize}
Beyond the simultaneous restoration of the chiral flavor and axial $U(1)_{A}$ symmetries, the persistent splitting between the axial-vector and tensor channels ($m_{AV} - m_{Tt}$) carries deep implications for the underlying nature of the Quark-Gluon Plasma (QGP). The lack of $SU(2)_{CS}$ chiral-spin symmetry restoration up to 204.5 MeV indicates that the thermal medium in this temperature window is far from a non-interacting ideal gas of quarks and gluons. Rather, it implies that strong non-perturbative chromomagnetic interactions remain unscreened and dynamically active, characterizing the QGP as a strongly coupled fluid.While these preliminary results provide a rigorous non-perturbative baseline for understanding the phase structure of the QCD vacuum, they also define clear directions for future investigations. Our current simulations on a 48³ × 16 lattice are inherently subject to finite-volume and finite-lattice-spacing effects. To fully map the thermodynamic limit and perform a robust continuum limit extrapolation ($a \rightarrow 0$), future studies will require finer lattice spacings and lighter, strictly physical quark masses. Scaling up this computationally intensive MDWF framework will rely heavily on next-generation exascale high-performance computing architectures, such as the Fugaku supercomputer, to map the critical scaling behaviors with unprecedented statistical accuracy.
